\section{Quinta iteración de la Fase 3: Lanzar las rutas}
\label{sec:LanzarRutas}

	
	Con el generador de aleatorios verificado esta iteración le da
	su primer uso real: elegir rutas. Hasta ahora el motor de estados
	enciende segmentos con una SECUENCIA de prueba escrita en el código;
	esa secuencia desaparece. A partir de aquí, lo que se enciende y se
	apaga en el cristal son \textbf{rutas}: cadenas de segmentos que la
	autoría dibujó con los enlaces a \texttt{siguiente}, y que empiezan en
	los segmentos de tipo \textbf{inicio}, que el diseñador también numeró
	(campo \texttt{ruta} del GWI). El motor elige qué ruta lanza de forma 
	aleatória; el elemento la recorre; cuando termina, desaparece y al tick 
	siguiente nace otra.
	
	Será la primera vez en que algo definido en el cartucho, y no constantes de
	prueba, definen lo que se ve en pantalla. 
	
	De momento solo se presentan los segmentos en las rutas, no se va a verificar las Colisiones con su enlace a vida o muerte, aquí las rutas nacen y continúna hasta el final de su recorrido, lo que permite tanto probar esta características del motor de estados, como visualmente apreciar que las secuencias definidas son correctas.	
	
	\subsection{Qué es una ruta en el modelo}
	
	Una ruta es la cadena que ya viaja en el GWI desde el diseñador:
	un segmento \textbf{inicio} (marcado con un número de \texttt{ruta}) y la
	cadena segmentos definidos por los enlaces \texttt{siguiente} que salen de él. El inicio de la	ruta es quien da el número: el proceso lanza una ruta por su número, se localiza el segmento de inicio correspondiente, y se lanza la cadena.
	
	El \emph{elemento} que recorre la ruta no es una entidad nueva:
	es ``el segmento que está encendido ahora en esa ruta''. Como todo
	en el modelo, el estado vive en la lista de segmentos: recorrer una
	ruta es apagar el segmento actual y encender el que marca su enlace a
	\texttt{siguiente} — el mismo mecanismo que ya usa el jugador, solo
	que el jugador se mueve con botones y la ruta avanza con el reloj.
	
	\subsection{Decisiones actuales}
	
	\begin{enumerate}	
		\item \textbf{¿Cuándo se lanza una nueva ruta?} Se lanzará \emph{en el mismo tick en que muere la anterior}, de forma
		que las rutas van encadenadas sin huecos. Mientras una ruta vive, no
		nace ninguna otra; el tick en que la ruta muere, el PRNG elige la
		siguiente y la arranca. La partida siempre tiene una ruta en marcha
		(salvo antes del primer tick).
		
		Esta política de lanzamiento es para esta fase de pruebas; en una fase
		posterior, la política, las aceleraciones, lanzar varias rutas a
		la vez estilo Game B, hacer pausas más cortas, etc. viajarán en el
		cartucho como parámetros. 
		
		\item \textbf{¿Qué hace el elemento al llegar al final de la ruta
			sin colisión?} \emph{Desaparece}: la ruta muere al final.
		El último segmento de la cadena se apaga y no se sigue
		ningún enlace. Los enlaces \texttt{salvado}/\texttt{muerte}
		se tratarán cuando se aborde el tema de las colisiones, que será la que
		decida cuál de los dos rutas se toma.
		
		\item El PRNG es el ya definido, sin cambios: el motor pide
		«dame una secuencia» al director con parámetro
		(\texttt{aleatorio(maxRutas)}) y el director responde con la
		secuencia determinista generada desde la semilla. Misma semilla, misma
		partida.
	\end{enumerate}
	
	\subsection{El tick del motor, con rutas}
	
	El orden de las acciones del tick no cambia (entrada, apagar
	banderas de cambiado, avance); lo que cambia es el paso de avance:
	la SECUENCIA de prueba se sustituye por el avance de la ruta y el
	lanzamiento de la nueva:
	
	\begin{itemize}
		\item En el barrido solo participan los segmentos de tipo juego (1) e
		inicio (2). Los segmentos del jugador (0) y los
		contadores de vida/muerte (3 y 4) jamás son barridos aunque estén
		encendidos. Para verificarlo se crea una función \texttt{esDeRuta} que nos lo indicará.
		
		\item Antes de tocar nada se toma una FOTO de los
		segmentos de ruta visibles: cada uno avanza exactamente un paso,
		aunque su \texttt{siguiente} también estuviera visible. Sin esta salvaguardia una cadena podría avanzar dos pasos en un mismo tick.
		
		\item \textbf{Avance:} si hay ruta viva, su segmento actual se
		apaga y se enciende su \texttt{siguiente} (ambos quedan
		marcados como cambiados). Si el actual no tiene
		\texttt{siguiente} (el final de la cadena), la ruta MUERE:
		el segmento se apaga y la ruta deja de existir. Esto ocurre ahora por no controlar colisiones, pero pueden existir cadenas que hagan exactamente esto para lanzar cosas por pantalla que no intervienen en el juego, como en nuestro caso de ejemplo, pueden aparecer tiburones nadando bajo la barca que no intervienen en el juego, solo son un adorno en la pantalla.
		
		\item \textbf{Lanzamiento:} si no hay ruta viva, el motor pide un
		número de ruta al PRNG del director, localiza el segmento de
		inicio de esa ruta y lo enciende: la ruta nace. Nace en el
		mismo tick en que se nota la ausencia, así que la partida
		nunca muestra un tick ``vacío'' (salvo el arranque).
		
		\item \textbf{Arranque de la partida} (\texttt{iniciar}): se
		apagan todas las banderas y NO hay ruta viva: el primer
		tick lanza la primera. El jugador sigue buscándose y
		encendiéndose como hasta ahora: la ruta y el
		jugador conviven en pantalla sin conocerse.
	\end{itemize}
	
	El motor no añade ningún estado propio: si una ruta vive, es porque
	su segmento actual tiene la bandera \texttt{visible} encendida; si
	no vive ninguna, es que ningún segmento de ruta la tiene. Como
	siempre, la lista de segmentos es el estado: activo y visible son
	lo mismo a efectos de presentación. (Hay una función de diagnóstico,
	\texttt{rutasVivas()}, que cuenta elementos de ruta visibles, pero
	no es estado: es una lectura de la lista).
	
	El mecanismo no limita el número de rutas simultáneas: si algún
	día se lanzan varias a la vez, cada una avanza por su propio
	\texttt{siguiente} en el mismo barrido. Es la política de
	lanzamiento (una sola viva) lo que impone el encadenado, no el
	barrido.
	
	\subsection{Detalles que la simulación debe cazar}
	
	\begin{itemize}
		\item La primera ruta nace en el primer tick tras
		\texttt{iniciar()} (no en el propio \texttt{iniciar}).
		
		\item La ruta recorre la CADENA completa: la simulación debe
		verificar la secuencia exacta de segmentos encendidos en
		pantalla (a nivel de píxel o de banderas), no solo ``algo
		se enciende''.
		
		\item Al morir la ruta, el segmento final se apaga SIN RASTRO
		(el repintado completo del motor gráfico ya lo garantiza)
		y al tick siguiente nace la nueva.
		
		\item Reproducibilidad: con la misma semilla, dos partidas
		eligen el MISMO orden de rutas (el PRNG ya está verificado;
		aquí se comprueba que el consumo de números es estable).
		
		\item Regresiones: el jugador sigue moviéndose con los botones
		por sus enlaces mientras las rutas corren; la cola de
		pulsaciones, la marca de primer jugador y el motor gráfico
		no se tocan.
	\end{itemize}
	
	\subsection{Pruebas}
	
	Poniendo en marcha el juego, ya se aprecia que de forma aleatoria se van lanzando las tres rutas hasta el final, y siempre que se reinicia el juego comienza la misma secuencia.
	
\clearpage